% ============================================================
\section{Introduction}
\label{sec:introduction}
% ============================================================

\subsection{Liquid Biopsy and the Clinical Need for Rapid Variant Detection}

The analysis of cell-free DNA (cfDNA) circulating in peripheral blood has
emerged as one of the most consequential advances in precision oncology over the
past decade~\cite{wan2017liquid}.  Unlike solid-tissue biopsies, which are
invasive, single-site snapshots of tumour heterogeneity, liquid biopsies
provide minimally invasive access to the collective genetic landscape of all
tumour foci through the tumour-derived fraction of cfDNA, termed circulating
tumour DNA (\ctDNA{})~\cite{siravegna2017integrating}.  Clinical applications
span screening for early-stage cancer, treatment selection by detection of
actionable mutations, longitudinal monitoring of therapeutic response, and early
detection of relapse through minimal residual disease (MRD) surveillance.

The MRD monitoring context imposes the most stringent demands on sensitivity.
After curative-intent treatment, \ctDNA{} may persist at variant allele
frequencies (VAF) as low as $0.01\text{--}0.1\%$ of total cfDNA, and the
ability to detect this ultra-low-level signal provides early warning of
impending relapse, often weeks to months before clinical symptoms appear or
imaging findings are detectable~\cite{bettegowda2014detection,heitzer2019current}.
At VAF $= 0.1\%$ with $3000\times$ sequencing depth, a true somatic variant
will produce only $\sim$3 supporting reads per target position --- a signal
barely distinguishable from sequencing error rates of $0.1\text{--}1\%$.

The standard bioinformatics pipeline for somatic variant detection --- alignment
of reads to a reference genome, followed by variant calling with tools such as
GATK Mutect2~\cite{cibulskis2013mutect} or Strelka2~\cite{saunders2012strelka}
--- achieves excellent sensitivity for germline variants and for somatic
mutations at moderate VAF ($>5\%$), but its performance degrades substantially
in the ultra-low-VAF regime characteristic of MRD monitoring.  Furthermore, the
alignment-variant-calling pipeline is computationally intensive (1--2 hours for
a typical targeted panel), incompatible with same-day clinical reporting.

\subsection{The K-mer Variant Detection Paradigm}

K-mer-based variant detection offers a complementary approach that sidesteps
the alignment step entirely~\cite{audoux2017km,audoux2020km}.  Rather than
aligning reads to a reference and then identifying positions where aligned
bases differ, k-mer analysis pre-computes a count database of all distinct
substrings of length $k$ (k-mers) in the sequencing data, then \emph{walks}
through k-mer space from a known reference sequence to discover all variant
paths supported by the data.

The core insight is that a single nucleotide variant (SNV) changes exactly $k$
consecutive k-mers relative to the reference; an insertion of length $L$ creates
$L + k - 1$ novel k-mers; a deletion removes junction k-mers and creates $k - 1$
new spanning k-mers.  If the count database is queried for all possible single-base
extensions of each reference k-mer, branching points in the resulting graph
reveal the locations of variants, and the counts along each path quantify their
allele frequency.

This approach has several compelling properties for the liquid biopsy setting.
First, k-mer counting is entirely reference-free in its counting phase: all
k-mers present in the sequencing data are counted without any prior alignment.
This eliminates reference bias and avoids the loss of reads that align poorly
due to large insertions or tandem duplications, which are particularly common
clinical targets (e.g., FLT3-ITD in AML).  Second, k-mer counts are directly
proportional to molecular abundance, providing a natural linear model for
variant allele frequency estimation.  Third, with modern k-mer counters such as
\jf{}~\cite{marcais2011fast}, the pre-computation phase is extremely fast
($\sim$1.5 minutes for a typical targeted sequencing run), and individual k-mer
lookups are $O(1)$ from the resulting hash table, enabling rapid processing of
many targets in parallel.

\subsection{The km/kmtools/kam Ecosystem and Its Limitations}

The \km{} tool (IRIC Bioinformatics, Audoux et al.~\cite{audoux2017km,audoux2020km})
implements the k-mer walking algorithm described above and has been validated for
targeted variant detection in haematological malignancies, particularly for
detection of FLT3-ITD and NPM1 insertion mutations in AML.  \km{} is a
Python tool that interfaces with \jf{} through its Python bindings and produces
tabular output describing all variant paths detected for each target sequence.

The \kmtools{} companion package and the \kam{} pipeline
(Trentzsch~\cite{kamtool}) extend \km{} with parallelism (distributing targets
across threads), filtering (matching detected variants against known patient
mutations), merging of multi-run outputs, statistical summaries, and
visualisation.  The \kam{} Docker container integrates the entire workflow from
raw FASTQ reads through deduplicated k-mer counting (via HUMID~\cite{pockrandt2020humid})
to filtered variant calls.

Validation of the \kam{} pipeline on a 10-patient cohort of liquid biopsy MRD
samples sequenced with UMI-based duplex sequencing revealed several significant
limitations.  SNV sensitivity reached only 77\%, substantially below the
90--95\% achieved by alignment-based callers on the same samples.  INDEL
sensitivity was dramatically lower at 38\% overall and near-zero ($\sim$3\%) for
insertions longer than 7~bp, primarily due to fundamental constraints of the
k-mer length relative to insertion size.  The practical VAF detection threshold
was approximately 0.1\%, limiting the approach's utility for the deepest MRD
monitoring scenarios.

Beyond sensitivity, the \km{}/\kmtools{}/\kam{} ecosystem has architectural
limitations.  It is implemented in Python with C++ \jf{} bindings, requiring
separate processes for deduplication, counting, detection, and filtering.
Parallelism is achieved by pre-distributing target files across subdirectories
(\textit{refolder}) rather than through native thread-level parallelism.  The
DFS walking algorithm uses fixed extension thresholds that cannot adapt to
sample-specific coverage depth or error rate.  Variant calls carry no
statistical quality annotation: there are no p-values, confidence intervals on
rVAF, strand bias scores, or Phred-scaled quality metrics, making the output
incompatible with standard VCF-based downstream analysis.

\subsection{kmerdet: Goals and Contributions}

\kmerdet{} is a Rust reimplementation and significant algorithmic extension of
the \km{}/\kmtools{}/\kam{} ecosystem, designed to address each of the
limitations described above while maintaining full backward compatibility with
existing \km{} outputs and target catalogs.

The primary contributions of \kmerdet{} are:

\begin{enumerate}

  \item \textbf{Single-binary pipeline}: \kmerdet{} replaces the
    multi-process \kam{} pipeline (\textit{HUMID} + \textit{jellyfish} +
    \textit{multiseqex} + \textit{refolder} + \textit{kmtools chunk} +
    \textit{kmtools filter}) with a unified binary providing subcommands
    \texttt{detect}, \texttt{filter}, \texttt{merge}, \texttt{stats},
    \texttt{plot}, \texttt{coverage}, and \texttt{run}, eliminating process
    spawn overhead, intermediate file I/O, and external dependency management.

  \item \textbf{Native parallelism}: Target-level parallelism is implemented
    using Rayon's work-stealing scheduler~\cite{rayon}, automatically
    distributing work across all available cores without the file-organisation
    overhead of \textit{refolder} + \textit{kmtools chunk}.

  \item \textbf{Bidirectional k-mer walking}: The DFS walker extends k-mers
    both forward (rightward) and backward (leftward) from every reference
    k-mer, improving detection of variants near target boundaries and deletions
    that span reference anchor regions.

  \item \textbf{Coverage-adaptive extension thresholds}: Extension thresholds
    are computed from a per-sample Poisson noise model calibrated against the
    local reference k-mer coverage, replacing the fixed \texttt{ratio} and
    \texttt{count} parameters that cannot adapt to coverage variation between
    samples and targets.

  \item \textbf{Multi-k consensus detection}: Running the walking algorithm at
    multiple k values ($k = 21, 31, 41$ or per-target adaptive selection)
    with weighted consensus voting substantially improves INDEL sensitivity by
    providing more spanning k-mers for large insertions at shorter k.

  \item \textbf{Statistical quality annotation}: Each variant call is
    annotated with a binomial p-value (combined across k-mer positions using
    Fisher's method~\cite{fisher1925statistical}), a Phred-scaled QUAL score,
    and a 95\% bootstrap confidence interval on rVAF, enabling standard
    VCF-compatible filtering.

  \item \textbf{Panel-of-normals filtering}: A built-in panel-of-normals (PoN)
    subcommand builds and applies a database of recurrent artifacts observed in
    healthy controls, suppressing systematic sequencing errors, CHIP variants,
    and common germline polymorphisms.

  \item \textbf{Graph pruning}: Post-walking tip removal, bubble collapsing,
    and dead-end pruning reduce graph complexity and eliminate short error-
    derived branches before pathfinding, improving both precision and speed.

  \item \textbf{VCF 4.3 output}: Full VCF 4.3~\cite{vcfspec} output with
    standard INFO/FORMAT fields (DP, AF, QUAL, FS, TLOD) alongside the
    existing TSV format, enabling direct integration with downstream analysis
    pipelines.

\end{enumerate}

The guiding principle throughout is \emph{correctness before cleverness}: every
algorithmic deviation from \km{}'s behaviour is intentional, documented, and
validated against the reference implementation.  On identical inputs, \kmerdet{}
produces output that matches \km{}'s TSV format field-for-field.

\subsection{Paper Structure}

The remainder of this paper is organised as follows.
Section~\ref{sec:background} reviews the biological context of liquid biopsy
ctDNA monitoring, relevant genomics methods, and the statistical challenges of
ultra-low-VAF variant detection.
Section~\ref{sec:methods} presents the complete \kmerdet{} algorithm, from
k-mer encoding through walking, graph construction, pathfinding, variant
classification, and NNLS quantification, including all algorithmic improvements
over \km{}.
Section~\ref{sec:implementation} describes the Rust implementation, key design
decisions, and the type-driven architecture that makes the codebase both safe
and extensible.
Section~\ref{sec:results} presents benchmark results comparing \kmerdet{} with
the \kam{} pipeline on the validation cohort, with placeholder structure for
planned benchmarking experiments.
Section~\ref{sec:discussion} discusses the clinical implications of the
improvements, remaining limitations, and future research directions.
Section~\ref{sec:conclusion} concludes.
